% * =========================================================================
% * 后期拟完成的研究工作及进度安排
% * =========================================================================
\section{后期拟完成的研究工作及进度安排}

本课题所定核心目标已基本完成，后期工作以成果深化、优化机理验证、方法论凝练与论文撰写为主，具体包括：
\begin{semiQuotList}
    \item 持续完善学术论文，根据IPDPS 2026的审稿意见进行修改与补充实验，提升工作的完整性与说服力。
    \item 深化对已有四项系统级优化（CFWP、NMA、EPM、TPS）生效机理的验证：在原有端到端吞吐与消融实验之外，补充直接刻画各优化点作用机制的细粒度实验，详见第一点说明。
    \item 在已有系统优化方法论的基础上规划新的创新点，探索将移位异或纠删码扩展应用于大语言模型推理的KV Cache（键值缓存）存储场景（如Mooncake），详见第二点说明。
    \item 系统整理研究成果，撰写硕士毕业论文并准备答辩。
\end{semiQuotList}

\subsection{第一点：已有系统优化的机理验证与实验深化}

现有学术论文主要通过端到端吞吐量与消融（Ablation）实验，从结果层面验证了四项优化（缓存友好窗口划分CFWP、相邻合并访问NMA、统一解码抽象EPM、遍历模式选择TPS）带来的整体性能提升。
后期拟在此基础上补充更直接刻画各优化点\emph{作用机制}的细粒度实验，使“为何有效”与“有效多少”都得到充分支撑，具体规划如下。

\textbf{缓存命中率验证（对应CFWP）。}
缓存友好窗口划分的核心收益在于提升缓存局部性，但现有实验仅以吞吐间接反映。
后期拟借助\texttt{perf}或Cachegrind等工具直接采集L1/L2/LLC（末级缓存）的缓存命中率与缺失率（Cache Miss Rate），对比开启/关闭窗口划分及不同窗口大小$W$下的指标变化，直接证明缓存局部性的改善，并据此印证窗口大小取值的合理性。

\textbf{内存访问量与写放大验证（对应NMA）。}
相邻合并访问的理论收益是将校验块写入访问量约减半。
后期拟通过访存计数（如\texttt{perf}内存事件或插桩统计）测量编码过程中对数据块、校验块的实际读写次数，验证其与理论复杂度$\Theta((k+\tfrac{km}{2})l)$的吻合程度，量化写放大（Write Amplification）的降低幅度。

\textbf{冗余访存消除验证（对应EPM）。}
统一解码抽象通过将数据块解码与校验块修复融合，省去独立的修复阶段。
后期拟对比统一解码流程与“解码＋独立修复”两阶段流程的总访存次数与遍历轮数，量化被消除的冗余内存访问，直接解释EPM带来约一倍解码性能提升的来源。

\textbf{单块丢失快路径验证（对应TPS）。}
遍历模式选择针对最常见的单块丢失场景采用纵向单遍异或快路径。
后期拟分别统计单块丢失与多块丢失场景下的访存模式与每码字访问次数，验证快路径“每码字恰好访问一次”的特性，并测量流式（Streaming）指令绕过缓存所带来的额外收益。

\textbf{实验维度扩展。}
配合上述机理验证，进一步补充不同块大小、不同$(k,m)$配置与更多丢失模式下的细粒度数据，并对$(10,4)$等不利配置下的性能瓶颈做针对性分析与优化探索。

\subsection{第二点：面向KV Cache存储的移位异或纠删码（新增创新点规划）}

针对毕业设计在学术论文基础上的扩展创新点，本课题拟将移位异或纠删码引入大语言模型推理的KV Cache（Key-Value Cache，键值缓存）存储场景，初步规划如下。

\textbf{问题定位。}
大语言模型推理过程中，预填充（Prefill）阶段产生的KV Cache规模巨大（如LLaMA3-70B在128k tokens下的KV Cache可达约40\,GB），需借助分布式缓存池（以Mooncake为代表，采用以KV Cache为中心的解耦架构）跨请求、跨实例复用。
此类缓存池规模大、节点故障频繁，现有方案多依赖多副本（Replication）保证可用性，导致内存利用率偏低、等效缓存容量与命中率受限。

\textbf{机会点。}
相比多副本，纠删码能以更低的冗余开销同时提供存储节省与故障容错；然而KV Cache的重建处于推理的延迟关键路径上，传统Reed-Solomon码因含有限域乘法、编解码开销较大，难以满足低延迟要求。
而移位异或码仅依赖移位与异或操作，配合本课题已完成的系统级优化（FastTT），具备极高的编解码吞吐，恰好为纠删码进入KV Cache关键路径提供了可行性。

\textbf{研究内容。}
拟将Shift-XOR/FastTT的高吞吐编解码引入KV Cache缓存池，将纠删码重新定位为“性能导向的冗余”：以同一套编码机制，既降低缓存冗余开销以提升等效容量，又在节点故障或缓存驱逐时快速就地重建KV块。

\textbf{预期收益。}
一方面，用纠删码替代多副本可提升等效缓存容量，进而提高缓存命中率、减少昂贵的预填充重算；
另一方面，在节点故障时通过快速移位异或解码就地重建丢失的KV块，避免触发GPU重算，从而维持首Token延迟（Time To First Token，TTFT）等服务等级目标（Service Level Objective，SLO）。

\textbf{开放问题与计划。}
后期将进一步研究：KV块粒度与编码条带（Stripe）的对齐方式；解码重建延迟与预填充重算延迟之间的折中边界；纠删编码与缓存池哈希索引、驱逐（Eviction）策略的协同；以及冗余编码采用写时同步或后台异步以避免拖慢写入关键路径等问题。

具体工作进度安排为：
\begin{closeItemize}
    \item 2026年3月 -- 2026年4月：完善论文与补充实验，凝练移位异或纠删码系统优化方法论，规划毕业设计的新增创新点。
    \item 2026年4月 -- 2026年5月：开展新创新点的初步研究与验证，扩展实验，整理全部实验数据与结论。
    \item 2026年5月 -- 2026年6月：撰写硕士毕业论文，完善图表与方法论总结，准备毕业答辩。
\end{closeItemize}
